% ============================================================
%  LECSEG — A1 defense poster
%  Uses tikzposter. Single page, A1 portrait.
% ============================================================

\documentclass[a1paper, portrait, 25pt]{tikzposter}
\usepackage{graphicx}
\usepackage{booktabs}
\usepackage[utf8]{inputenc}
\usepackage{xcolor}

\usetheme{Default}
\usecolorstyle{Russia}

\newcommand{\lecseg}{\textsc{LecSeg}}
\newcommand{\dataset}{\textsc{LecSeg-30}}
\newcommand{\todo}[1]{\textcolor{red}{\textbf{[#1]}}}

\title{Hierarchical Multimodal Lecture-Video Topic Segmentation}
\author{Author One, Author Two, Author Three, Author Four, Author Five
        \quad | \quad Supervisor: Dr.\ Supervisor Name}
\institute{BRAC University, Department of Computer Science and Engineering}

\begin{document}
\maketitle

% ----- Row 1 -----
\begin{columns}
  \column{0.5}
  \block{Why lecture-video segmentation?}{
    \todo{2--3 sentences on the access/searchability problem and one image
    of an unindexed long lecture.}
  }
  \column{0.5}
  \block{Our contributions}{
    \begin{itemize}
      \item Open hierarchical multimodal pipeline.
      \item Reliability-weighted fusion.
      \item Local-LLM refinement (no closed APIs).
      \item \dataset{}: 30 videos, 5 domains, dual annotation.
      \item 5-metric evaluation with bootstrap CIs.
    \end{itemize}
  }
\end{columns}

% ----- Row 2: pipeline -----
\block{The \lecseg{} pipeline}{
  \todo{Wide pipeline diagram (5 boxes). Caption underneath.}
}

% ----- Row 3: results -----
\begin{columns}
  \column{0.5}
  \block{Main results}{
    \todo{Bar chart of $P_k$ across 8 methods.}
  }
  \column{0.5}
  \block{Ablations}{
    \todo{4-bar ablation summary.}
  }
\end{columns}

% ----- Row 4: take-aways -----
\block{Take-aways}{
  \begin{itemize}
    \item \lecseg{} reduces $P_k$ by \todo{X.X}\,\% relative ($p<\todo{0.0X}$).
    \item Hierarchical output adds \todo{Y.Y}\,\% improvement on the new HWD metric.
    \item Local LLM matches closed-API alternatives at \todo{Z}\,\% the cost.
  \end{itemize}
}

% ----- Footer -----
\block*{Resources}{
  Code: \texttt{github.com/yourorg/lecseg} \quad
  Dataset: \texttt{zenodo.org/records/...} \quad
  Model: \texttt{huggingface.co/yourorg/lecseg-llama-8b}
}

\end{document}
